\section{Related Work}
\label{sec:related_work}

Our work resides at the intersection of parameter-space model merging, dynamic activation-space blending, and hardware-aware conditional inference.

\subsection{Parameter-Space Model Merging}
Model merging has emerged as a promising paradigm to consolidate multiple specialized neural networks into a single cohesive model without undergoing expensive retraining~\citep{wortsman2022model,choshen2022fusing}. Early methods focused on simple parameter-space averaging (e.g., Model Soups~\citep{wortsman2022model}). While successful for models fine-tuned on the same downstream task, naive parameter averaging fails when merging models fine-tuned on diverse, heterogeneous tasks due to severe weight interference~\citep{yadav2023ties}. To mitigate this interference, Task Arithmetic~\citep{ilharco2022editing} introduces linear arithmetic on task-specific weight vectors. More recently, TIES-Merging~\citep{yadav2023ties} resolves weight conflicts by identifying and keeping only sign-consistent dominant parameters, and DARE~\citep{yu2023language} prunes small weight changes before merging. 

While parameter-space merging maintains the exact inference latency of the base model, it is fundamentally static. It collapses all specialized adapters into a single static set of weights, causing performance degradation under high task diversity—a phenomenon known as ``heterogeneity collapse''~\citep{yadav2023ties}. 

\subsection{Dynamic Activation-Space Blending}
To overcome the limitations of static merging, several recent works advocate for dynamic, sample-wise activation-space blending. Instead of merging weights offline, these techniques run parallel expert adapters and dynamically blend their activations at runtime based on the input query's features. SABLE~\citep{sable2025} employs task classification heads to route sample activations dynamically. SPS-ZCA~\citep{spszca2025} resolves the temporal paradox of routing by pre-computing task centroids in the early layers (Layer 3) and applying zero-shot centroid alignment (ZCA). ChemMerge~\citep{chemmerge2026} models continuous state tracking of routing weights using chemical kinetics to avoid layer-wise activation spikes. 

Although activation-space blending preserves specialized task performance and avoids heterogeneity collapse, it has a severe deployment drawback: \textbf{uncontrolled computational scale}. Running $K$ parallel adapter pathways across all layers increases both the floating-point operations (FLOPs) and DRAM weight transfer bandwidth, causing thermal throttling and high latency on embedded devices~\citep{shen2024collaborative}. Q-SPS~\citep{qsps2026} attempts to mitigate this by quantizing activations and experts and using a static coefficient threshold to skip expert execution. However, Q-SPS relies on static, hardcoded thresholds and cannot adapt to dynamic hardware events (like low battery or OS-level resource reallocation) in real time.

\subsection{Low-Power Edge Serving and TinyML}
Optimizing deep learning for resource-constrained edge devices (TinyML) is a highly active field~\citep{warden2020tinyml}. Traditional approaches include model pruning~\citep{han2015deep}, quantization~\citep{gholami2021survey}, and knowledge distillation. Dynamic inference techniques, such as early exiting (e.g., BranchyNet~\citep{teerapittayanon2016branchynet}) and slimmable networks~\citep{tang2020slimmable}, adapt a model's depth or width to reduce latency on-the-fly. Additionally, Sparse Mixture of Experts (MoE)~\citep{shazeer2017outrageously,fedus2022switch} routes tokens to a subset of experts to save compute. 

However, existing dynamic edge-serving models generally require extensive end-to-end retraining or fine-tuning, and their gating structures are hardwired into the network architecture. In contrast, RB-TopM is a completely training-free, plug-and-play framework. It operates on top of pre-trained, unmodified LoRA adapters and base models, utilizing a system-level hardware control loop ($C_{\text{budget}}$) to dynamically regulate both the expert capacity and pruning aggression in microseconds at runtime. This unique capability makes RB-TopM highly practical and immediately deployable in real-world volatile edge environments.
